% GENERATED FILE. Edit canonical lessons, then run npm run generate:books.
% canonical-source-hash: fnv1a64:883693497156aa54
% canonical-lessons: ML-C566-kalimannu, ML-C566-cessu, ML-C566-cittukali, ML-C566-tottappani, ML-C566-dhyanikkuka

\chapter{Clay, Cards and Gardening}
\label{ch:ml-clay-cards-and-gardening}
\hlchaptermodality{566}

\begin{chapteropening}
\textbf{By the end of this chapter:} \emph{I can say clay, chess, a card game, gardening and to meditate.}

Complete the last lesson of chapter 566: I can say clay, chess, a card game, gardening and to meditate.
\end{chapteropening}

\section[\texorpdfstring{\ml{കളിമണ്ണ്} (kaḷimaṇṇŭ)}{കളിമണ്ണ് (kaḷimaṇṇŭ)}]{\ml{കളിമണ്ണ്} (kaḷimaṇṇŭ) — clay}

\label{lesson:ML-C566-kalimannu}



\begin{quote}
Before the new one: say the Malayalam for a writer, then the Malayalam for to produce a film.
\end{quote}

\subsection*{You'll want to know: \ml{കളിമണ്ണ്}}
\textbf{\ml{കളിമണ്ണ്}} — \emph{kaḷimaṇṇŭ} — \textquotedblleft{}clay\textquotedblright{}.

\begin{quote}
\ml{കുട്ടികൾ} \ml{കളിമണ്ണ്} \ml{കൊണ്ട്} \ml{ചെറിയ} \ml{പാത്രങ്ങൾ} \ml{ഉണ്ടാക്കി}.
\end{quote}

\begin{quote}
\emph{kuṭṭikaḷ kaḷimaṇṇŭ koṇṭŭ ceṟiya pātraṅṅaḷ uṇṭākki.}
\end{quote}

\begin{quote}
The children made little pots out of clay.
\end{quote}

The potter's earth; its second part is \ml{മണ്ണ്}, earth.

\begin{cousinweb}
\textbf{\ml{കളിമണ്ണ്}} (\emph{kaḷimaṇṇŭ}) is native Dravidian: DEDR 1381 gives the Malayalam compound itself as potter's clay, built on \ml{മണ്ണ്}, \textquotedblleft{}earth\textquotedblright{} (DEDR 4666).
\end{cousinweb}

\begin{grammarlens}[title={What you've built}]
One more word for describing what you lived and felt.
\end{grammarlens}

\subsection*{Your turn}
\begin{itemize}
\raggedright
  \item \emph{Say it:} \emph{kaḷimaṇṇŭ}
  \item \emph{Say it:} \emph{kaḷimaṇṇŭ}, once more
  \item \emph{Say it:} \emph{kuṭṭikaḷ kaḷimaṇṇŭ koṇṭŭ ceṟiya pātraṅṅaḷ uṇṭākki.} — The children made little pots out of clay.
  \item \emph{Recall:} say the Malayalam for a writer, then the Malayalam for to produce a film, then say \emph{kaḷimaṇṇŭ} again
\end{itemize}

\begin{tcolorbox}[breakable,colback=teal!4,colframe=teal!35!black,title={Before you move on}]
What does \ml{കളിമണ്ണ്} mean? (\textquotedblleft{}Clay\textquotedblright{}.) Say it once more.
\end{tcolorbox}
\section[\texorpdfstring{\ml{ചെസ്സ്} (cessŭ)}{ചെസ്സ് (cessŭ)}]{\ml{ചെസ്സ്} (cessŭ) — chess}

\label{lesson:ML-C566-cessu}



\begin{quote}
Before the new one: say the Malayalam for to produce a film, then the Malayalam for clay.
\end{quote}

\subsection*{You'll want to know: \ml{ചെസ്സ്}}
\textbf{\ml{ചെസ്സ്}} — \emph{cessŭ} — \textquotedblleft{}chess\textquotedblright{}.

\begin{quote}
\ml{ഒഴിവുസമയത്ത്} \ml{മുത്തച്ഛൻ} \ml{എന്നോടൊപ്പം} \ml{ചെസ്സ്} \ml{കളിക്കും}.
\end{quote}

\begin{quote}
\emph{oḻivusamayattŭ muttacchan ennōṭoppaṁ cessŭ kaḷikkuṁ.}
\end{quote}

\begin{quote}
In his free time grandfather plays chess with me.
\end{quote}

One plays it with \ml{കളിക്കുക}.

\begin{cousinweb}
\textbf{\ml{ചെസ്സ്}} (\emph{cessŭ}) was borrowed from English \emph{chess}.
\end{cousinweb}

\begin{grammarlens}[title={What you've built}]
One more word for describing what you lived and felt.
\end{grammarlens}

\subsection*{Your turn}
\begin{itemize}
\raggedright
  \item \emph{Say it:} \emph{cessŭ}
  \item \emph{Say it:} \emph{cessŭ}, once more
  \item \emph{Say it:} \emph{oḻivusamayattŭ muttacchan ennōṭoppaṁ cessŭ kaḷikkuṁ.} — In his free time grandfather plays chess with me.
  \item \emph{Recall:} say the Malayalam for to produce a film, then the Malayalam for clay, then say \emph{cessŭ} again
\end{itemize}

\begin{tcolorbox}[breakable,colback=teal!4,colframe=teal!35!black,title={Before you move on}]
What does \ml{ചെസ്സ്} mean? (\textquotedblleft{}Chess\textquotedblright{}.) Say it once more.
\end{tcolorbox}
\section[\texorpdfstring{\ml{ചീട്ടുകളി} (cīṭṭukaḷi)}{ചീട്ടുകളി (cīṭṭukaḷi)}]{\ml{ചീട്ടുകളി} (cīṭṭukaḷi) — a card game; cards}

\label{lesson:ML-C566-cittukali}



\begin{quote}
Before the new one: say the Malayalam for clay, then the Malayalam for chess.
\end{quote}

\subsection*{You'll want to know: \ml{ചീട്ടുകളി}}
\textbf{\ml{ചീട്ടുകളി}} — \emph{cīṭṭukaḷi} — \textquotedblleft{}a card game; cards\textquotedblright{}.

\begin{quote}
\ml{അവധിക്ക്} \ml{ഞങ്ങൾ} \ml{എല്ലാവരും} \ml{ചീട്ടുകളിയിൽ} \ml{ചേർന്നു}.
\end{quote}

\begin{quote}
\emph{avadhikkŭ ñaṅṅaḷ ellāvaruṁ cīṭṭukaḷiyil cērnnu.}
\end{quote}

\begin{quote}
In the holidays we all joined in a card game.
\end{quote}

\ml{ചീട്ട്} a playing card, also a slip of paper, + \ml{കളി} a game.

\begin{cousinweb}
\textbf{\ml{ചീട്ടുകളി}} (\emph{cīṭṭukaḷi}) is a Malayalam compound built from \ml{ചീട്ട്}, \textquotedblleft{}a card, a slip\textquotedblright{}, whose origin is not settled, and \ml{കളി}, \textquotedblleft{}play, a game\textquotedblright{}, which is native (DEDR 1374).
\end{cousinweb}

\begin{grammarlens}[title={What you've built}]
One more word for describing what you lived and felt.
\end{grammarlens}

\subsection*{Your turn}
\begin{itemize}
\raggedright
  \item \emph{Say it:} \emph{cīṭṭukaḷi}
  \item \emph{Say it:} \emph{cīṭṭukaḷi}, once more
  \item \emph{Say it:} \emph{avadhikkŭ ñaṅṅaḷ ellāvaruṁ cīṭṭukaḷiyil cērnnu.} — In the holidays we all joined in a card game.
  \item \emph{Recall:} say the Malayalam for clay, then the Malayalam for chess, then say \emph{cīṭṭukaḷi} again
\end{itemize}

\begin{tcolorbox}[breakable,colback=teal!4,colframe=teal!35!black,title={Before you move on}]
What does \ml{ചീട്ടുകളി} mean? (\textquotedblleft{}A card game; cards\textquotedblright{}.) Say it once more.
\end{tcolorbox}
\section[\texorpdfstring{\ml{തോട്ടപ്പണി} (tōṭṭappaṇi)}{തോട്ടപ്പണി (tōṭṭappaṇi)}]{\ml{തോട്ടപ്പണി} (tōṭṭappaṇi) — gardening}

\label{lesson:ML-C566-tottappani}



\begin{quote}
Before the new one: say the Malayalam for chess, then the Malayalam for a card game.
\end{quote}

\subsection*{You'll want to know: \ml{തോട്ടപ്പണി}}
\textbf{\ml{തോട്ടപ്പണി}} — \emph{tōṭṭappaṇi} — \textquotedblleft{}gardening\textquotedblright{}.

\begin{quote}
\ml{വിരമിച്ച} \ml{ശേഷം} \ml{അച്ഛന്റെ} \ml{ഹോബി} \ml{തോട്ടപ്പണിയാണ്}.
\end{quote}

\begin{quote}
\emph{viramicca śēṣaṁ acchanṟe hōbi tōṭṭappaṇiyāṇŭ.}
\end{quote}

\begin{quote}
Since he retired, Father's hobby is gardening.
\end{quote}

\ml{തോട്ടം} a garden + \ml{പണി} work.

\begin{cousinweb}
\textbf{\ml{തോട്ടപ്പണി}} (\emph{tōṭṭappaṇi}) is native Dravidian, a compound of \ml{തോട്ടം}, \textquotedblleft{}a garden\textquotedblright{} (DEDR 3549), and \ml{പണി}, \textquotedblleft{}work, labour\textquotedblright{} (DEDR 3884).
\end{cousinweb}

\begin{grammarlens}[title={What you've built}]
One more word for describing what you lived and felt.
\end{grammarlens}

\subsection*{Your turn}
\begin{itemize}
\raggedright
  \item \emph{Say it:} \emph{tōṭṭappaṇi}
  \item \emph{Say it:} \emph{tōṭṭappaṇi}, once more
  \item \emph{Say it:} \emph{viramicca śēṣaṁ acchanṟe hōbi tōṭṭappaṇiyāṇŭ.} — Since he retired, Father's hobby is gardening.
  \item \emph{Recall:} say the Malayalam for chess, then the Malayalam for a card game, then say \emph{tōṭṭappaṇi} again
\end{itemize}

\begin{tcolorbox}[breakable,colback=teal!4,colframe=teal!35!black,title={Before you move on}]
What does \ml{തോട്ടപ്പണി} mean? (\textquotedblleft{}Gardening\textquotedblright{}.) Say it once more.
\end{tcolorbox}
\section[\texorpdfstring{\ml{ധ്യാനിക്കുക} (dhyānikkuka)}{ധ്യാനിക്കുക (dhyānikkuka)}]{\ml{ധ്യാനിക്കുക} (dhyānikkuka) — to meditate}

\label{lesson:ML-C566-dhyanikkuka}



\begin{quote}
Before the new one: say the Malayalam for a card game, then the Malayalam for gardening.
\end{quote}

\subsection*{You'll want to know: \ml{ധ്യാനിക്കുക}}
\textbf{\ml{ധ്യാനിക്കുക}} — \emph{dhyānikkuka} — \textquotedblleft{}to meditate\textquotedblright{}.

\begin{quote}
\ml{അമ്മ} \ml{എല്ലാ} \ml{ദിവസവും} \ml{രാവിലെ} \ml{പത്ത്} \ml{മിനിറ്റ്} \ml{ധ്യാനിക്കും}.
\end{quote}

\begin{quote}
\emph{amma ellā divasavuṁ rāvile pattŭ miniṟṟŭ dhyānikkuṁ.}
\end{quote}

\begin{quote}
Every morning Mother meditates for ten minutes.
\end{quote}

Past: \ml{ധ്യാനിച്ചു}, \emph{dhyāniccu}.

The noun is \ml{ധ്യാനം}, meditation.

\begin{cousinweb}
\textbf{\ml{ധ്യാനിക്കുക}} (\emph{dhyānikkuka}) is built from Sanskrit \emph{dhyāna}, \textquotedblleft{}meditation\textquotedblright{} (CDIAL 6889), with the Malayalam verb ending -\ml{ിക്കുക}.
\end{cousinweb}

\begin{grammarlens}[title={What you've built}]
One more word for describing what you lived and felt.
\end{grammarlens}

\subsection*{Your turn}
\begin{itemize}
\raggedright
  \item \emph{Say it:} \emph{dhyānikkuka}
  \item \emph{Say it:} \emph{dhyānikkuka}, once more
  \item \emph{Say it:} \emph{amma ellā divasavuṁ rāvile pattŭ miniṟṟŭ dhyānikkuṁ.} — Every morning Mother meditates for ten minutes.
  \item \emph{Recall:} say the Malayalam for a card game, then the Malayalam for gardening, then say \emph{dhyānikkuka} again
\end{itemize}

\begin{tcolorbox}[breakable,colback=teal!4,colframe=teal!35!black,title={Before you move on}]
What does \ml{ധ്യാനിക്കുക} mean? (\textquotedblleft{}To meditate\textquotedblright{}.) Say it once more.
\end{tcolorbox}
